\section[Tests first]{Tests first}
\sectionframe{07}{Tests first}
  {How does a test written before the code keep an agent on track?}
  {\faCheckDouble}

% ---------------------------------------------------------------- the loop
\begin{frame}[source={Beck (2002), Test-Driven Development: By Example}]{Test-driven development: red, green, refactor}
  \begin{tikzpicture}[x=1mm,y=1mm,
      phase/.style={box,text width=28mm,minimum height=13mm,font=\footnotesize,line width=0.8pt},
      gain/.style={inner sep=0pt,font=\footnotesize,text=udsdark,align=left},
      tick/.style={inner sep=0pt,font=\footnotesize,anchor=base east}]
    \useasboundingbox (0,-1.36) rectangle (144,-55);
    \node[phase,draw=udsred,fill=redtint] (red) at (72,-8)
      {{\color{udsred}\bfseries\small RED}\\[0.4mm] write a failing test for one behaviour};
    \node[phase,draw=testgreen,fill=greentint] (green) at (96,-40)
      {{\color{testgreen}\bfseries\small GREEN}\\[0.4mm] write just enough code to pass it};
    \node[phase,draw=udsblue,fill=bluetint] (ref) at (48,-40)
      {{\color{udsblue}\bfseries\small REFACTOR}\\[0.4mm] tidy the code; every test stays green};
    \draw[flow] (red.east) to[bend left=35] (green.north);
    \draw[flow] (green.west) -- (ref.east);
    \draw[flow] (ref.north) to[bend left=35] (red.west);
    \node[lbl,align=center] at (72,-27) {one behaviour\\per cycle};
    % what each step gives us; each tick sits on its gain's first baseline (with a text width,
    % a node's base is its first line)
    \node[gain,anchor=west] (g1) at (95,-5.2) {the requirement, written first};
    \node[gain,anchor=west] (g2) at (95,-10.8) {proof that the test can fail};
    \node[gain,anchor=west,text width=20mm] (g3) at (118.4,-40) {small steps,\\easy to review};
    \node[gain,anchor=west,text width=22mm] (g4) at (8,-40) {a regression net\\for later changes};
    \foreach \g/\c in {g1/udsred,g2/udsred,g3/testgreen,g4/udsblue}
      \node[tick,text=\c] at ([xshift=-1.4mm]\g.base west) {\faCheck};
    \node[gain,anchor=base] (g5) at (74,-53.4) {tests call clear interfaces: the seams};
    \node[tick,text=udsdark] at ([xshift=-1.4mm]g5.base west) {\faCheck};
  \end{tikzpicture}
  \note{\textbf{33:40--34:40}\par
    Test-driven development in Kent Beck's formulation is a short loop.\par
    Red: write a test for one behaviour and watch it fail.\par
    Green: write just enough code to make it pass.\par
    Refactor: tidy the code while all tests stay green.\par
    Then the next behaviour.\par
    What we get: the requirement written down as a test before the code exists, proof that each
      test can fail, small steps, a regression net, and pressure towards clear interfaces, the
      seams where tests observe behaviour.}
\end{frame}

% ------------------------------------------------------- why for an agent
\begin{frame}[source={Anthropic (2026), Claude Code best practices · Mathews \& Nagappan (2024), ASE ’24}]{Why tests first help an agent}
  \begin{tikzpicture}[x=1mm,y=1mm,
      node2/.style={box,minimum height=9mm,inner sep=1.6mm},
      item/.style={anchor=base west,inner sep=0pt,font=\footnotesize,text=udsdark},
      head/.style={anchor=base west,inner sep=0pt,font=\footnotesize\bfseries}]
    \useasboundingbox (0,0) rectangle (144,-54);
    % the agent edits the code and runs the tests
    \node[node2,artifact,font=\ttfamily\footnotesize\bfseries,minimum width=20mm] (src) at (12,-8) {src/};
    \node[node2,agent,font=\footnotesize\bfseries,minimum width=26mm] (ag) at (62,-8) {\faRobot\enspace agent};
    \node[node2,artifact,draw=testgreen,fill=greentint,font=\ttfamily\footnotesize\bfseries,minimum width=20mm]
      (tst) at (114,-8) {tests/};
    \draw[flow] (ag.west) -- node[lbl,above] {edits} (src.east);
    \draw[flow] (ag.east) -- node[lbl,above] {runs} (tst.west);
    \node[anchor=south,inner sep=0pt,font=\scriptsize,text=testgreen] at ([yshift=1mm]tst.north)
      {\faLock\enspace approved and committed first};
    % pass or fail comes back: the signal
    \draw[flow,draw=testgreen,line width=0.8pt] (tst.south) to[out=-90,in=-90,looseness=0.55]
      node[pos=0.5,below,inner sep=0.8mm,font=\scriptsize\bfseries,text=testgreen] {pass or fail} (ag.south);
    % what the signal gives the agent
    \node[head,text=testgreen] at (0,-28.4) {what it gives the agent};
    \foreach \t [count=\i] in {a check it runs itself,{the intent, fixed first},one small goal per cycle,
        earlier decisions still enforced}
      \node[item] at (0,-28.4-\i*4.7) {\textcolor{testgreen}{\faCheck}\enspace\t};
    % how it gets gamed: aimed at the tests, and what catches it
    \node[head,text=udsred] (gamed) at (70,-28.4) {how it gets gamed};
    \foreach \t [count=\i] in {tests edited or skipped,expected values copied from the code,
        a mock that replaces the logic,tests fitted after the code}
      \node[item] at (70,-28.4-\i*4.7) {\textcolor{udsred}{\faTimes}\enspace\t};
    \draw[udsred,line width=0.6pt,dashed,-{Stealth[length=1.7mm,width=1.3mm]}]
      ([xshift=1.5mm]gamed.east) to[out=0,in=-90] (121,-12.6);
    \node[item,text=udsblue] at (70,-53)
      {caught by: test diffs reviewed first · CI runs every test};
  \end{tikzpicture}
  \note{\textbf{34:40--35:40}\par
    Anthropic's guidance for Claude Code puts it plainly: without a check it can run, an agent
      stops when the work looks done.\par
    A failing test turns done into a signal the agent reads and iterates on.\par
    Written first, from the spec, the test records intent instead of mirroring the
      implementation, and the suite keeps earlier decisions enforced in later sessions.\par
    The failure modes are real: agents weaken, skip or edit tests, recompute expected values the
      way the code does, or replace the logic under test with a mock, a stand-in that returns
      whatever the test expects.\par
    Hence our guardrails.\par
    A benchmark study found that giving the model the tests with the task solved more problems;
      that is evidence for code generation, not for research code.}
\end{frame}

% ------------------------------------------------- /implement drives /tdd
\begin{frame}[source={Pocock (2026), mattpocock/skills v1.3.1: implement, tdd · aihero.dev/skills}]{How \texttt{/implement} uses \texttt{/tdd}}
  \begin{columns}[T,onlytextwidth]
    \begin{column}{0.48\textwidth}
      \begin{tikzpicture}[x=1mm,y=1mm,
          stp/.style={agent,text width=48mm,minimum height=6.6mm,align=left,font=\scriptsize,inner sep=1mm},
          num/.style={font=\scriptsize\bfseries,text=udsblue,anchor=east,inner sep=0.8mm}]
        \node[anchor=north west,inner sep=0pt,font=\footnotesize] at (0,0)
          {\fn{/implement}: one ticket, one fresh session};
        \node[stp,anchor=north west] (s1) at (6,-5)  {read the ticket};
        \node[stp,anchor=north west,draw=testgreen,fill=greentint,line width=0.8pt,minimum height=9mm] (s2) at (6,-14)
          {\fn{/tdd} at the agreed seams: one failing test, then just enough code to pass it};
        \node[stp,anchor=north west] (s3) at (6,-26) {typecheck; run single test files};
        \node[stp,anchor=north west] (s4) at (6,-35) {run the full suite once};
        \node[stp,anchor=north west] (s5) at (6,-44) {\fn{/code-review}, then commit};
        \foreach \i in {1,...,5} \node[num] at (s\i.west) {\i};
        \foreach \i [evaluate=\i as \j using int(\i+1)] in {1,...,4} \draw[flow] (s\i.south) -- (s\j.north);
        \draw[flow,draw=testgreen] ([yshift=-2.2mm]s2.east) -- ++(3.5,0) -- ++(0,4.4)
          node[midway,right,lbl,text=testgreen] {repeat} -- ++(-3.5,0);
      \end{tikzpicture}
    \end{column}
    \begin{column}{0.49\textwidth}
      \begin{gblock}[fontupper=\footnotesize]{Rules of the \texttt{/tdd} loop}
        \begin{itemize}\setlength\itemsep{0.5mm}
          \item seams agreed with you first
          \item one failing test, then minimal code
          \item one vertical slice per cycle
          \item tests only through public interfaces
          \item mocks only at system boundaries
          \item refactoring left to \fn{/code-review}
        \end{itemize}
      \end{gblock}
      {\footnotesize Known gaps: code written before its test; ticket left open; the run reviews its own
       work.\par}
    \end{column}
  \end{columns}
  \note{\textbf{35:40--36:50}\par
    In Pocock's skills, implement is the build step and you invoke it yourself, one ticket per
      fresh session.\par
    It reads the ticket, drives the tdd skill at the agreed seams one red-green slice at a time,
      typechecks and runs single test files as it goes, runs the full suite once at the end,
      calls code-review and commits.\par
    The tdd skill is a set of rules: seams agreed first, red before green, one vertical slice
      per cycle, tests through public interfaces, mocks only at system boundaries.\par
    The loop has no refactor step, because agents almost never do it; code-review
      proposes refactoring instead.\par
    The author's documentation lists the gaps: the agent sometimes writes code first, implement
      does not close the ticket, and the review inside the same run reviews its own work, so a
      fresh session makes a better reviewer.}
\end{frame}

% ------------------------------------------------------- running example
\begin{frame}[fragile]{Red and green on the specimen split}
  \begin{columns}[T,onlytextwidth]
    \begin{column}{0.57\textwidth}
      % the tints run unbroken into the rules, and stop at the table's right edge
      \footnotesize\renewcommand\arraystretch{1.3}\setlength\tabcolsep{2.5pt}
      \setlength\aboverulesep{0pt}\setlength\belowrulesep{0pt}\setlength\extrarowheight{0.6pt}
      \begin{tabularx}{\linewidth}{@{}>{\leavevmode\bfseries\color{udsblue}}p{2.5mm}L{29mm}>{\columncolor{redtint}}L{22mm}>{\columncolor{greentint}[\tabcolsep][0pt]}X@{}}
        \toprule
        & \bfseries Test: the behaviour & \textbf{\textcolor{udsred}{Red: failure}} &
          \textbf{\textcolor{testgreen}{Green: fix}}\\
        \midrule
        1 & each specimen lands in exactly one fold & function does not exist & one fold per specimen ID\\
        2 & folds do not depend on row order & folds follow the row order & sort IDs, shuffle with the seed\\
        3 & missing \fn{specimen\_id} raises \fn{ValueError} & pandas raises \fn{KeyError} & check the columns first\\
        4 & too few specimens raise \fn{ValueError} & empty folds pass silently & check the count\\
        \bottomrule
      \end{tabularx}
    \end{column}
    \begin{column}{0.4\textwidth}
\begin{terminal}[basicstyle=\ttfamily\fontsize{7}{8.5}\selectfont]{terminal, cycle 1}
$ pytest tests/data -q
E   ImportError: cannot import name
    'specimen_folds'
ERROR tests/data/test_splitting.py
§\textcolor{termred}{1 error}§

# write the smallest specimen_folds()
$ pytest tests/data -q
.
§\textcolor{termgreen}{1 passed}§
\end{terminal}
    \end{column}
  \end{columns}
  \note{\textbf{36:50--38:00}\par
    Ticket T1 under the tdd loop.\par
    Cycle one: the test says every specimen lands in exactly one fold; it fails because the
      function does not exist; the smallest change assigns one fold per specimen ID.\par
    Cycle two: shuffling the manifest rows must not change the folds; the first version follows
      the row order, so this fails; sorting the IDs and shuffling them with the seed fixes
      it.\par
    Cycles three and four pin down the error cases, as the coding standards ask.\par
    The seams come from the spec, and the expected values come from a small hand-written fixture
      rather than from recomputing what the code does.\par
    The terminal output is abridged.}
\end{frame}
